\label{sec:methods}
\subsection{Semi-analytic gold standard}
Dividing $dS/dt$ by $dR/dt$ in \eqref{eq:sir} gives $dS/dR = -(\beta/\gamma N) S$, so
\begin{equation}
  S(t) = S_0 \exp\!\Big(-\tfrac{\beta}{\gamma N}\big(R(t)-R_0\big)\Big), \label{eq:phase}
\end{equation}
an exact, solver-free relation between $S$ and $R$ (we write $R_0$ here for the initial value of the
$R$ compartment, not to be confused with the reproduction number). Substituting \eqref{eq:phase} into
$dR/dt = \gamma I = \gamma(N-R-S)$ gives a \emph{scalar} separable ODE, so
\begin{equation}
  t(R) = \int_0^{R} \frac{dr}{\gamma\big(N - r - S_0 e^{-\frac{\beta}{\gamma N} r}\big)}. \label{eq:tofr}
\end{equation}
We evaluate \eqref{eq:tofr} by adaptive Simpson quadrature and invert it by a safeguarded
Newton method (the standard \texttt{rtsafe} scheme) to obtain $R(t)$, then recover $S$ from
\eqref{eq:phase} and $I = N-S-R$. This reference trajectory is accurate to machine precision using
no ODE solver at all, and every accuracy claim in this paper is measured against it. We additionally
use \eqref{eq:phase} to define the \emph{phase invariant}
$Q(t) = \ln S(t) + \frac{\beta}{\gamma N} R(t)$, constant along any true trajectory; unlike the
trivial mass-conservation check $S+I+R=N$ (which every explicit Runge--Kutta method satisfies to
roundoff at \emph{any} step size, since the right-hand side of \eqref{eq:sir} sums to zero
identically), drift in $Q(t)$ scales with the solver's true convergence order and is a genuine
accuracy diagnostic.

\subsection{Solvers}
We implement Forward Euler (order 1), Heun's method (order 2), classical RK4 (order 4), adaptive
Dormand--Prince RK45 with a PI step controller (order 5(4), used as the ``truth'' solver for model
variants with no closed form), and two implicit methods -- backward Euler and the trapezoidal rule
-- whose nonlinear stage equations are solved by Newton's method with the linear system at each
iteration factored by our own partial-pivoting LU decomposition.

\subsection{Estimation}
Given noisy observations $\mathrm{obs}_k$ at times $t_k$, we minimize the weighted least-squares cost
\begin{equation}
  J(\theta) = \sum_k w_k \big(\mathrm{obs}_k - h(y(t_k;\theta))\big)^2 \label{eq:cost}
\end{equation}
by grid search, golden-section coordinate descent, Nelder--Mead, Gauss--Newton, and
Levenberg--Marquardt, all implemented from scratch. The Jacobian required by the last two methods is
obtained by integrating the \emph{forward sensitivity equations}
$\dot s_j = J_y f \, s_j + \partial f/\partial\theta_j$ alongside the state, using the same solver
under test -- so solver error enters the Jacobian exactly as it enters the trajectory, which is the
mechanism behind our central claim, rather than by finite differences.

\subsection{Uncertainty quantification}
We compute four independent confidence regions for $(\hat\beta,\hat\gamma)$: an asymptotic ellipse
from $\mathrm{Cov}(\hat\theta) \approx \hat\sigma^2 (J^\top J)^{-1}$; a residual bootstrap; a Monte
Carlo ensemble over fresh noise realizations from the known ground truth (the only route that
measures true bias); and an adaptive Metropolis--Hastings sampler with four chains, reporting the
Gelman--Rubin $\hat R$ and effective sample size.

\subsection{The crossover experiment}
For a fixed $(\beta,\gamma)$, sampling interval, and fitting solver/step size $h$, we generate
noise-free synthetic data from the gold standard and fit with that solver: any deviation from the
true parameters is attributable purely to the solver's truncation error, giving a deterministic
$\mathrm{bias}(h)$. We then add noise at level $\sigma$ and refit over many replicates, giving
$\mathrm{std}(\sigma, h)$. The crossover $\sigma^*(h)$ is where $\mathrm{std}(\sigma^*, h) =
|\mathrm{bias}(h)|$.
